\documentclass{assignment}

\usepackage{hyperref}
\hypersetup{
    colorlinks=true,
    linkcolor=blue,
    urlcolor=blue
}

\begin{document}
\assignmentTitle{Pham Xuan
Hoang}{2640009}{assets/usth\_logo.png}{Information Systems}{Labwork 1}

\section*{Section 1: Installing MySQL}
I use CachyOS, a Linux distribution based on Arch Linux, so MySQL is
not available in the default package manager. I can install it
through the AUR (Arch User Repository) but I opted to use MariaDB, an
open source alternative (forked from MySQL).

On Arch, MariaDB is currently at version 13.x, which is newer than
the version that the test database is validated against (up to 12.x).
Therefore, I chose the lts (long-term support) version which is
currently at 12.x. I installed it
through \texttt{shelly} (a wrapper for \texttt{pacman}). And there
are some additional setup that is required to setup the database on
the system
(\href{https://wiki.archlinux.org/title/MariaDB\#Installation}{MariaDB
Installation on Arch wiki}).

\lstinputlisting[language=cli]{assets/install\_mariadb.txt}

Then I can start the service and check if the server is up and
listening on \texttt{0.0.0.0}.

\lstinputlisting[language=cli]{assets/start\_mariadb.txt}

\section*{Section 2: Set up account}
With the database up and running, I need to set up access on local
machine next. The first run required \texttt{sudo mariadb} because
only \texttt{root} and \texttt{mysql} account have access to the database.

I created a user with the same name as my local account so I
do not have to type \texttt{sudo} every time when I access the
database from my local machine.

\lstinputlisting[language=sql]{assets/create_main_user_mariadb.sql}

I also created another account for \texttt{tabularis}, an open source
GUI database client that I use.

\lstinputlisting[language=sql]{assets/create_client_user_mariadb.sql}

And we can check whether our account is set up correctly:

\lstinputlisting[language=sql]{assets/user\_check.txt}

\section*{Section 3: Import test data into MariaDB}
With the set up finished, I can init the database with
\texttt{test\_db} data and verify the correctness by running the SHA2 test.

\lstinputlisting[language=cli]{assets/import\_data.txt}

And then we can use \texttt{mariadb} client to check the table within
the \texttt{employees} database.

\lstinputlisting[language=sql]{assets/check\_data.txt}
\end{document}
